\section{Nonduality}

Originally published in June, 2005 for an open seminar.

\begin{quotex}
In all of us there dwells a secret marvelous power of freeing ourselves from the changes of time, of withdrawing to our secret selves away from external things, and of so discovering to ourselves the eternal in us in the form of unchangeability. This presentation of ourselves to ourselves is the most truly personal experience upon which depends everything that we know of the supra-sensual world. This presentation shows us for the first time what real existence is, whilst all else only appears to be. It differs from every presentation of the sense in its perfect freedom, whilst all other presentations are bound, being overweighted by the burden of the object. Still there exists for those who have not this perfect freedom of the inner sense some approach to it, experiences approaching it from which they may gain some faint idea of it. ... This intellectual presentation occurs when we cease to be our own object, when, withdrawing into ourselves, the perceiving self merges in the self-perceived. At that moment we annihilate time and duration of time; we are no longer in time, but time, or rather eternity itself, is in us. The external world is no longer an object for us, but is lost in us.
\flright{\textsc{Friedrich Schelling}, \emph{Philosophical Letters upon Dogmatism and Criticism}}
\end{quotex}


\paragraph{Dual Vision}

In dualistic vision there is the so-called ``real'' world ``out there, right now'' and a more or less exact copy in consciousness. Unfortunately, no one can explain how the mind creates a copy of the world. Moreover, all the qualities that make up consciousness, such as color, smell, etc., are not properties of the objects themselves.

Philosophers then dispute whether the ``I'' is part of, or distinct from, the space-time continuum.

\paragraph{Nondual Vision}

In nondual vision, objects arise in consciousness, emanating from the Noumenon. The Witness watches this arising, yet is not the cause of it, nor is it among the objects that it witnesses. Therefore, it can never know itself as an object.


\paragraph{Superimposition}

In superimposition, the dualistic view is superimposed on the primordial nonduality, leading to a mistaken understanding of reality.

The empirical ego, itself simply another object in space-time, then posits itself as the true self, clouding the unchanging and transcendent witness.

\textbf{Wei Wu Wei} is prodding us to make this shift in awareness from centered in the empirical ego as object, to the transcendent witness as subject. This is not so much a task to be done, as a non-action to let happen.

{Reference: \emph{The Open Secret}, Wei Wu Wei}


\flrightit{Posted on 2011-04-18 by Cologero}

